%قبل از شروع مسئله ترجیحاً چیزی ننویسید. اما اگر خواستید بنویسید، در قالب لیستی (فقرات و شمارش) بنویسید.
%برای اضافه کردن نمره، بعدی از شروع‌مسئله، [x] را اضافه کنید و به جای x نمره مورد نظر را قرار دهید. برای زیرمسئله‌ها هم به همین صورت است.
%\شروع‌مسئله[10]
\شروع‌مسئله [12]
\\
فرض کنید یک جدول درهم سازی با $m$ خانه داریم که خانه‌های آن را با اندیس‌های $\{0,1,...,m-1\}$ شماره گذاری کرده‌ایم.برای یافتن عنصر با کلید $k$ در این جدول مطابق الگوریتم زیر عمل می‌کنیم. \\
\شروع{شمارش}
 \فقره شمارنده‌های $h(k) \to x$ و $0 \to y$ را مقدار دهی کن.
 \فقره اگر کلید $k$ در خانه $x$ام جدول بود یا این خانه خالی بود جست‌و‌جو را متوقف کن.
 \فقره $x$ و $y$ را به شکل $m$ $mod$ $y+1 \to y$ و $m$ $mod$ $(x+y) \to x$ به‌روز کن و به مرحله‌ی ۲ برگرد.
\پایان{شمارش}

\شروع{زیرمسائل}
\زیرمسئله[4]
ثابت کنید الگوریتم فوق یک نوع وارسی درجه ۲ است.

\زیرمسئله[8]
حال فرض کنید $m = 2^d$ که $d$ یک عدد صحیح مثبت است. ثایت کنید الگوریتم فوق در بدترین حالت تمام خانه‌های جدول را بررسی می‌کند.
\پایان{زیرمسائل}

\solution{
آ)
در مرحله
$i$اُم
اجرای الگوریتم فوق خانه‌ی وارسی شده را می‌توان به شکل زیر به‌دست آورد.
$$h(k,i) = x(i) = (h(k) + 1 + 2 + ... + i + 1) \mod {m} = (h(k) + \frac{i(i+1)}{2}) \mod{m}$$
بنابراین این الگوریتم، یک الگوریتم وارسی درجه دو است.

ب)
حالتی که جدول پر است را در نظر بگیرید. کافیست ثابت کنیم تابع
$h(k,i)$
با دامنه
$0 \leq i \leq 2^d-1$
تابعی یک به یک است. (چرا؟)
بنابراین فرض کنید
$i$ و $j$
به گونه‌ای هستند که
$h(k, i) = h(k, j)$.
در این صورت
$$h(k, i) = h(k, j) \implies \frac{i(i+1)}{2} \equiv \frac{j(j+1)}{2} \pmod{2^d} \implies (i-j)(i+j+1) \equiv 0 \pmod{2^{d+1}}$$
از طرفی
$i \leq 2^d - 1$
و
$j \leq 2^d - 1$
بنابراین
$i + j + 1 \leq 2^{d+1} - 1$.
بنابراین
$i + j - 1$
بر
$2^{d+1}$
بخش‌پذیر نیست. بنابراین
$$i - j \equiv 0 \pmod{2^{d+1}}$$
که با توجه به اینکه
$i$ و $j$
هر دو کمتر از
$2^d$
هستند از عبارت فوق می‌توان نتیجه گرفت
$i = j$.
بنابراین خواسته ما اثبات شد.

}

%\متن‌سیاه{سوال عملی اول:}\\
%\متن‌سیاه{دراپ‌تست}\\

%\متن‌سیاه{ورودی:}\\
%در خط اول ورودی، ابتدا k و سپس m داده می‌شود.\\
%\متن‌سیاه{خروجی:}\\
%در تنها سطر خروجی، تعداد آیفون‌های مورد نیاز می‌آید. اگر در هیچ صورتی چنین کاری ممکن نبود ۱- چاپ می‌شود.\\
%\متن‌سیاه{محدودیت‌ها:}\\
%$1\leq k\leq 10^4$

%$m\leq10^9$
%\\
%\textbf{\متن‌سیاه{ورودی و خروجی نمونه}\\
%% تست اول
%\begin{latin}
% \begin{tabularx}{\textwidth}{|X | X|}
% \hline
% \textbf{stdin} &
% \textbf{stdout}\\
% \hline
% 100 100&
% 1\\
% \hline
% 
% \end{tabularx}
%\end{latin}
%
%% تست دوم
%
%\begin{latin}
% \begin{tabularx}{\textwidth}{|X | X|}
% \hline
% \textbf{stdin} &
% \textbf{stdout}\\
% \hline
% 340 20&
% 3\\
% \hline
% 
% \end{tabularx}
%\end{latin}}
% صورت اصلی مسئله را در این‌جا بنویسید.



\solution{
}

%%%%%%%%%%%%%%%%%%%%%%%%%%%%%%%%%%%%%%%%%%%%%%%%%%%%%
% برخی دستورات مفید برای تمرین‌ها. %
%%%%%%%%%%%%%%%%%%%%%%%%%%%%%%%%%%%%%%%%%%%%%%%%%%%%%

% فرمول‌های نمونه:
%$f(x) = \Theta(x)$
%$T(x, y) \leq \log(x) + 2^y + \binom{2n}{n}$
% {\tt code\_function}


% لیست موارد بدون شماره‌گذاری:
%\شروع{فقرات}
% \فقره مورد اول 
% \شروع{فقرات}
% \فقره مورد اول داخلی 
% \شروع{فقرات}
% \فقره اولین مورد خیلی داخلی 
% \فقره دومین مورد خیلی داخلی 
% \پایان{فقرات}
% \فقره مورد دوم داخلی 
% \پایان{فقرات}
% \فقره مورد دوم 
%\پایان{فقرات}

% لیست موارد شماره‌گذاری شده:
%\شروع{شمارش}
% \فقره یک مورد 
% \فقره یک مورد دیگر 
% \فقره باز هم یک مورد دیگر
%\پایان{شمارش}

% برای نوشتن یک فرمول چند خطی که یک شماره بخوره. برای حالت‌های دیگر به سایت زیر مراجعه کنید:
% https://www.overleaf.com/learn/latex/Aligning_equations_with_amsmath
%\begin{align} 
% \begin{split}
% (x+y)^3 &= (x+y)^2(x+y) \\
% &= (x^2+2xy+y^2)(x+y) \\
% &= (x^3+2x^2y+xy^2) + (x^2y+2xy^2+y^3) \\
% &= x^3+3x^2y+3xy^2+y^3
% \end{split} 
%\end{align}

% نوشتن ماتریس و آرایه دوبعدی
%\begin{align}
% A = 
% \begin{bmatrix}
% A_{11} & A_{21} \\
% A_{21} & A_{22}
% \end{bmatrix}
%\end{align}

% اضافه کردن تصویر با شرح. برای اضافه کردن تصویر بدون شرح، فقط دستور درج‌تصویر را استفاده کنید.
%\شروع{شکل}[h!]
%\تنظیم‌ازوسط
% \درج‌تصویر[پهنا=0.5\پهنای‌متن]{img.jpg}
%\شرح{شرح تصویر}
%\برچسب{تصویر:۱}
%\پایان{شکل}
%\پایان{زیرمسائل}

%کد زیر را ببینید.
%%\lstinputlisting{source_filename.py}
%\begin{latin}
%\begin{lstlisting}
%# Import the modules
%import sys
%import random
%
%ans = True
%
%while ans:
%question = raw_input("Ask the magic 8 ball a question: (press enter to quit) ")
%
%answers = random.randint(1,8)
%
%if question == "":
%sys.exit()
%
%elif answers == 1:
%print "It is certain"
%
%elif answers == 2:
%print "Outlook good"
%
%elif answers == 3:
%print "You may rely on it"
%
%elif answers == 4:
%print "Ask again later"
%
%elif answers == 5:
%print "Concentrate and ask again"
%
%elif answers == 6:
%print "Reply hazy, try again"
%
%elif answers == 7:
%print "My reply is no"
%
%elif answers == 8:
%print "My sources say no"
%\end{lstlisting}
%\end{latin}